
\documentclass[10pt,aspectratio=169]{beamer}
\usepackage[utf8]{inputenc}
\usepackage[T1]{fontenc}
\usepackage[catalan]{babel}
\usepackage{amsmath,amssymb,booktabs,array,multicol}
\usepackage{tikz,pgfplots,xcolor,listings}
\usetikzlibrary{positioning,arrows.meta,calc,shapes.geometric}
\pgfplotsset{compat=1.18}
\usepgfplotslibrary{fillbetween}

\definecolor{UABGreen}{RGB}{0,103,71}
\definecolor{UABLight}{RGB}{224,241,236}
\definecolor{DeepBlue}{RGB}{20,65,110}
\definecolor{AccentOrange}{RGB}{230,126,34}
\definecolor{SoftGray}{RGB}{245,247,250}
\definecolor{CodeBack}{RGB}{248,248,248}
\definecolor{CodeFrame}{RGB}{210,215,220}

\usetheme{Madrid}
\usecolortheme[named=UABGreen]{structure}
\setbeamercolor{title}{fg=white,bg=UABGreen}
\setbeamercolor{frametitle}{fg=white,bg=UABGreen}
\setbeamercolor{block title}{fg=white,bg=DeepBlue}
\setbeamercolor{block body}{fg=black,bg=SoftGray}
\setbeamercolor{alerted text}{fg=AccentOrange}
\setbeamertemplate{navigation symbols}{}
\setbeamertemplate{itemize item}{\color{UABGreen}$\blacktriangleright$}
\setbeamertemplate{itemize subitem}{\color{AccentOrange}$\bullet$}
\setbeamertemplate{enumerate item}{\color{UABGreen}\insertenumlabel.}
\setbeamertemplate{footline}{%
  \leavevmode%
  \hbox{%
  \begin{beamercolorbox}[wd=.72\paperwidth,ht=2.4ex,dp=1ex,leftskip=1.4em]{author in head/foot}%
    Estadística -- Grau en Enginyeria Informàtica
  \end{beamercolorbox}%
  \begin{beamercolorbox}[wd=.28\paperwidth,ht=2.4ex,dp=1ex,center]{date in head/foot}%
    \insertframenumber{} / \inserttotalframenumber
  \end{beamercolorbox}}%
  \vskip0pt%
}

\lstdefinelanguage{R}{%
  morekeywords={if,else,repeat,while,function,for,in,next,break,TRUE,FALSE,NULL,NA,NaN,Inf,
    library,mean,median,sd,var,summary,table,prop.table,xtabs,addmargins,round,
    ggplot,aes,geom_point,geom_line,geom_col,geom_histogram,geom_density,geom_smooth,
    geom_tile,facet_wrap,labs,theme_minimal,summarise,group_by,mutate,read_csv,filter,
    arrange,set.seed,rbinom,rgeom,rpois,runif,rexp,rnorm,dnorm,pnorm,qnorm,dt,pt,qt,
    dchisq,pchisq,qchisq,df,pf,qf,dbinom,pbinom,dpois,ppois,replicate,sample,rep,seq,
    tibble,data.frame,confint,t.test,var.test,chisq.test,binom.test,prop.test,lm,
    broom,glance,augment,tidy,future_map_dbl,plan,multisession,map_dbl,purrr,future,furrr},
  sensitive=true,
  morecomment=[l]\#,
  morestring=[b]",
  morestring=[b]'
}
\lstset{language=R,basicstyle=\ttfamily\scriptsize,keywordstyle=\color{DeepBlue}\bfseries,
  commentstyle=\color{gray},stringstyle=\color{AccentOrange},backgroundcolor=\color{CodeBack},
  frame=single,rulecolor=\color{CodeFrame},breaklines=true,columns=fullflexible,keepspaces=true,
  showstringspaces=false,
  literate={à}{{\`a}}1 {è}{{\`e}}1 {é}{{\'e}}1 {í}{{\'i}}1 {ï}{{\"i}}1 {ò}{{\`o}}1 {ó}{{\'o}}1 {ú}{{\'u}}1 {ü}{{\"u}}1 {ç}{{\c{c}}}1 {À}{{\`A}}1 {È}{{\`E}}1 {É}{{\'E}}1 {Í}{{\'I}}1 {Ò}{{\`O}}1 {Ó}{{\'O}}1 {Ú}{{\'U}}1}

\newcommand{\R}{\textsf{R}}
\newcommand{\RStudio}{\textsf{RStudio/Posit}}
\newcommand{\GitHub}{\textsf{GitHub}}
\newcommand{\important}[1]{\textcolor{UABGreen}{\textbf{#1}}}
\newcommand{\exemple}[1]{\textcolor{AccentOrange}{\textbf{#1}}}
\newcommand{\Prob}{\mathbb{P}}
\newcommand{\E}{\mathbb{E}}
\newcommand{\Var}{\operatorname{Var}}
\newcommand{\IC}{\operatorname{IC}}

\title[Probabilitat IV]{Probabilitat IV}
\author{Estadística -- Grau en Enginyeria Informàtica}
\date{Curs 2026--27}

\begin{document}
\begin{frame}[plain]
  \titlepage
  \vspace{-0.45cm}
  \begin{center}
  \small\textcolor{DeepBlue}{Estadística computacional amb \R, simulació i exemples d'enginyeria informàtica}
  \end{center}
\end{frame}


\begin{frame}[fragile]{Per què el TCL és central en estadística computacional?}
\begin{itemize}
  \item La mitjana de moltes contribucions independents tendeix a comportar-se com una normal.
  \item Això explica per què la normal apareix en errors de mesura, benchmarks repetits i estimadors.
  \item Permet quantificar incertesa sense conèixer exactament la distribució original.
\end{itemize}
\begin{block}{Idea clau}
Encara que les dades originals no siguin normals, la mitjana mostral pot ser aproximadament normal quan $n$ és prou gran.
\end{block}
\end{frame}

\begin{frame}[fragile]{Objectius del tema}
\begin{enumerate}
  \item Recordar independència, suma d'esperances i suma de variàncies.
  \item Formular el Teorema Central del Límit.
  \item Aplicar el TCL a sumes i mitjanes.
  \item Verificar-lo amb simulació Monte Carlo en \R.
  \item Connectar-lo amb estimació, intervals de confiança i benchmarking.
\end{enumerate}
\end{frame}

\begin{frame}[fragile]{Variables independents i agregació}
\begin{columns}[T]
\column{0.52\textwidth}
Si $X$ i $Y$ són independents:
\[
\E(X+Y)=\E(X)+\E(Y),
\]
\[
\Var(X+Y)=\Var(X)+\Var(Y).
\]
Per constants $\lambda$:
\[
\E(\lambda X)=\lambda\E(X),\quad
\Var(\lambda X)=\lambda^2\Var(X).
\]
\column{0.42\textwidth}
\begin{block}{Exemple}
La durada total de 45 components independents és la suma de 45 durades individuals.
\end{block}
\begin{alertblock}{Atenció}
La independència no és només una qüestió tècnica: si hi ha dependència, la variància agregada pot canviar molt.
\end{alertblock}
\end{columns}
\end{frame}

\begin{frame}[fragile]{Teorema Central del Límit}
\begin{block}{TCL clàssic}
Si $X_1,\ldots,X_n$ són independents i idènticament distribuïdes amb
\[
\E(X_i)=\mu,\qquad \Var(X_i)=\sigma^2,
\]
aleshores, per $n$ prou gran,
\[
\bar X=\frac{X_1+\cdots+X_n}{n}\approx N\left(\mu,\frac{\sigma}{\sqrt n}\right).
\]
\end{block}
\begin{itemize}
  \item La desviació típica de $\bar X$ és $\sigma/\sqrt n$.
  \item Augmentar $n$ redueix la incertesa de la mitjana.
\end{itemize}
\end{frame}

\begin{frame}[fragile]{Interpretació per a benchmarks}
\begin{columns}[T]
\column{0.56\textwidth}
Quan mesurem moltes vegades el temps d'execució d'un algoritme:
\begin{itemize}
  \item cada execució té soroll: càrrega del sistema, memòria cau, processos concurrents;
  \item la mitjana de moltes execucions és més estable;
  \item el TCL justifica aproximar la distribució de la mitjana amb una normal.
\end{itemize}
\column{0.38\textwidth}
\begin{block}{Regla pràctica}
No comparem algoritmes amb una sola execució. Repetim, resumim i quantifiquem incertesa.
\end{block}
\end{columns}
\end{frame}

\begin{frame}[fragile]{Exemple: vida total de components}
\begin{itemize}
  \item 45 components funcionen un darrere l'altre.
  \item Cada component té vida mitjana 50 h i desviació típica 4 h.
  \item La durada total és $X=X_1+\cdots+X_{45}$.
\end{itemize}
\[
\E(X)=45\cdot 50=2250,
\quad
\operatorname{sd}(X)=\sqrt{45\cdot 4^2}\approx 26.8.
\]
Pel TCL:
\[
X\approx N(2250,26.8).
\]
\begin{lstlisting}
mu <- 45 * 50
sigma <- sqrt(45 * 4^2)
1 - pnorm(2300, mean = mu, sd = sigma)
\end{lstlisting}
\end{frame}

\begin{frame}[fragile]{Simulació del TCL amb dades no normals}
\begin{lstlisting}
library(tidyverse)
set.seed(123)
B <- 10000
n <- 40
mitjanes <- replicate(B, mean(rexp(n, rate = 1)))

tibble(mitjanes) |>
  ggplot(aes(mitjanes)) +
  geom_histogram(aes(y = after_stat(density)), bins = 40) +
  geom_density(linewidth = 1) +
  theme_minimal()
\end{lstlisting}
\begin{block}{Observació}
Tot i que l'exponencial és asimètrica, les mitjanes simulades són molt més semblants a una normal.
\end{block}
\end{frame}

\begin{frame}[fragile]{Com canvia la distribució de la mitjana?}
\begin{columns}[T]
\column{0.46\textwidth}
\begin{itemize}
  \item Amb $n$ petit, la mitjana conserva part de la forma original.
  \item Amb $n$ gran, la forma s'aproxima a la normal.
  \item La variabilitat baixa com $1/\sqrt n$.
\end{itemize}
\column{0.48\textwidth}
\begin{lstlisting}
sim_mitjana <- function(n, B = 5000) {
  replicate(B, mean(rexp(n, rate = 1)))
}

sd(sim_mitjana(5))
sd(sim_mitjana(50))
\end{lstlisting}
\end{columns}
\end{frame}

\begin{frame}[fragile]{Aproximació normal de proporcions}
\begin{block}{Connexió amb la binomial}
Si $X\sim B(n,p)$, aleshores $\hat p=X/n$ és la mitjana de $n$ Bernoulli.
Pel TCL:
\[
\hat p \approx N\left(p,\sqrt{\frac{p(1-p)}{n}}\right).
\]
\end{block}
\begin{itemize}
  \item Exemple: taxa d'errors en 10.000 peticions.
  \item Base dels intervals de confiança i tests per proporcions.
\end{itemize}
\end{frame}

\begin{frame}[fragile]{Paral·lelitzar simulacions Monte Carlo}
\begin{lstlisting}
library(furrr)
plan(multisession, workers = 4)

simula <- function(n) mean(rexp(n, rate = 1))
mitjanes <- future_map_dbl(1:10000, ~ simula(100))

mean(mitjanes)
sd(mitjanes)
\end{lstlisting}
\begin{alertblock}{Per què té sentit paral·lelitzar?}
Cada repetició Monte Carlo és independent: és un cas natural de càlcul paral·lel.
\end{alertblock}
\end{frame}

\end{document}
